% Vietnamese translation for OI Public Library
% Source-File: IOI中国国家候选队论文/2007/day2/3.郭华阳《RMQ与LCA问题》.ppt
\documentclass[11pt]{article}
\usepackage{oipl-vn}

\newcommand{\Oh}{\mathcal{O}}
\newcommand{\RMQ}{\operatorname{RMQ}}
\newcommand{\LCA}{\operatorname{LCA}}
\newcommand{\pos}{\operatorname{pos}}
\newcommand{\ST}{\operatorname{ST}}

\title{Bản trình chiếu: RMQ và bài toán LCA}
\author{Guo Huayang}
\date{2007}

\begin{document}
\maketitle

\origtitle{RMQ và LCA}
\sourcepath{IOI China candidate team papers / 2007 / day 2 / Guo Huayang.ppt}

\section{Bố cục}

Tệp gốc là PowerPoint, không có PDF đi kèm trong thư mục nguồn. Phần chữ
trích được cho thấy deck có ba tuyến chính: bài toán RMQ, quan hệ hai chiều
giữa RMQ và LCA, rồi các ứng dụng và cách cài đặt như Kruskal, Tarjan,
sparse table và block RMQ.

\section{RMQ}

Với một mảng tĩnh \(A[1..n]\), bài toán \term{range minimum query} hỏi vị trí
hoặc giá trị nhỏ nhất trên đoạn \([i,j]\). Cấu trúc dữ liệu được đo bằng hai
tham số: thời gian tiền xử lý \(f(n)\) và thời gian trả lời truy vấn \(g(n)\).
Trong nhiều bài thi, dữ liệu không đổi nhưng số truy vấn lớn, nên bỏ thêm
tiền xử lý để đạt truy vấn hằng số là rất đáng giá.

Cách cơ sở là sparse table. Đặt
\[
  st[k][i]=\arg\min_{i\le t<i+2^k} A[t],
\]
hoặc lưu trực tiếp giá trị nhỏ nhất trên đoạn độ dài \(2^k\). Công thức dựng:
\[
  st[k+1][i]=\min(st[k][i],st[k][i+2^k])
\]
theo nghĩa so sánh giá trị trong \(A\). Tiền xử lý tốn \(\Oh(n\log n)\). Với
truy vấn \([l,r]\), lấy \(k=\lfloor\log_2(r-l+1)\rfloor\), so sánh hai đoạn
\([l,l+2^k-1]\) và \([r-2^k+1,r]\), vì chúng phủ toàn bộ đoạn hỏi. Truy vấn
tốn \(\Oh(1)\).

\section{LCA quy về RMQ}

Với cây có gốc, \(\LCA(u,v)\) là tổ tiên chung sâu nhất của \(u\) và \(v\).
Làm một Euler tour từ gốc: mỗi khi đi vào một đỉnh hoặc quay về từ con, ghi
đỉnh đó vào dãy \(E\), đồng thời ghi độ sâu tương ứng vào dãy \(B\). Gọi
\(\pos(u)\) là vị trí xuất hiện đầu tiên của \(u\) trong \(E\). Khi đó
\[
  \LCA(u,v)=E\bigl[\RMQ(B,\pos(u),\pos(v))\bigr],
\]
trong đó RMQ chọn vị trí có độ sâu nhỏ nhất trên đoạn. Trực giác là đoạn
Euler từ lần đầu gặp \(u\) đến lần đầu gặp \(v\) đi lên đúng tới tổ tiên chung
gần nhất rồi đi xuống; độ sâu nhỏ nhất trên đoạn chính là đỉnh ấy.

Vì Euler tour có độ dài \(2n-1\), ta có thể xử lý LCA bằng sparse table trong
\(\Oh(n\log n)\) tiền xử lý và \(\Oh(1)\) mỗi truy vấn. Nếu dùng RMQ tuyến
tính, LCA cũng đạt \(\Oh(n)\) tiền xử lý và \(\Oh(1)\) truy vấn.

\section{RMQ quy về LCA}

Chiều ngược lại cũng được slide nhắc tới. Từ mảng \(A\), dựng cây Cartesian:
phần tử nhỏ nhất của đoạn làm gốc, đoạn bên trái tạo cây con trái, đoạn bên
phải tạo cây con phải. Khi đó vị trí nhỏ nhất trên đoạn \([i,j]\) chính là
\(\LCA(i,j)\) trong cây Cartesian. Cây này có thể dựng bằng ngăn xếp đơn điệu
trong \(\Oh(n)\). Vì vậy một cấu trúc LCA tốt cũng cho lời giải RMQ tốt.

Quan hệ hai chiều này giải thích vì sao hai bài toán thường được học cùng
nhau: RMQ là mô hình mảng tĩnh, còn LCA là mô hình cây; Euler tour và cây
Cartesian là hai phép chuyển đổi giữa chúng.

\section{Tarjan LCA offline}

Nếu toàn bộ truy vấn LCA đã biết trước, thuật toán Tarjan dùng DFS và DSU để
trả lời trong thời gian gần tuyến tính. Tại đỉnh \(u\), ta tạo tập DSU cho
\(u\), duyệt từng con \(v\), xử lý xong cây con \(v\) thì hợp nhất \(v\) vào
\(u\) và đặt đại diện tổ tiên của tập mới là \(u\). Khi một truy vấn
\((u,v)\) được xét và đầu kia đã được thăm xong, đáp án là tổ tiên được gắn
với đại diện DSU của đầu kia.

Giả mã trên slide có dạng:
\[
\begin{array}{l}
\texttt{Tarjan\_DFS(u)}\\
\quad F(u)\leftarrow u\\
\quad \text{với mỗi truy vấn }(u,v)\text{ đã sẵn sàng, trả lời bằng }F(v)\\
\quad \text{với mỗi con }v\text{ của }u:\\
\quad\quad \texttt{Tarjan\_DFS(v)}\\
\quad\quad F(v)\leftarrow u.
\end{array}
\]
Trong cài đặt đầy đủ cần thêm màu thăm và thao tác \(\operatorname{find}\) của
DSU; độ phức tạp là \(\Oh((n+q)\alpha(n))\).

\section{Ứng dụng với Kruskal}

Deck nhắc một bài năm 2006 với các tham số \(N\le 1000\),
\(M,Q\le 100000\), \(D\le 5000\). Phần chữ trích được cho thấy lời giải dùng
Kruskal và truy vấn trên cây: các cạnh được xét theo thứ tự trọng số, DSU tạo
các lần hợp nhất, rồi những truy vấn \((u,v)\) được trả lời qua một cấu trúc
cây sinh bởi quá trình Kruskal.

Mô hình thường gặp là \term{cây Kruskal}. Khi hai thành phần được hợp nhất
bởi một cạnh trọng số \(w\), tạo một nút mới có trọng số \(w\) và nối nó với
gốc của hai thành phần. Trong cây này, LCA của hai đỉnh gốc ban đầu cho biết
thời điểm hoặc ngưỡng trọng số mà hai đỉnh lần đầu nằm trong cùng thành phần.
Vì vậy sau khi sắp cạnh tốn \(\Oh(M\log M)\), các truy vấn có thể được đưa về
LCA offline bằng Tarjan hoặc online bằng RMQ/sparse table.

Các công thức độ phức tạp trên slide có dạng
\(\Oh(M\log M + D N\alpha(N)+Q)\) hoặc \(\Oh(M\log M + Q + ND)\), tùy cách
tổ chức dữ liệu theo các lớp/trạng thái \(D\). Thông điệp chính là phần đồ
thị động theo ngưỡng được nén thành một cây, rồi LCA đảm nhiệm phần truy vấn.

\section{Block RMQ}

Sparse table đạt truy vấn \(\Oh(1)\) nhưng tiền xử lý \(\Oh(n\log n)\). Deck
còn nhắc cách đạt \(\Oh(n)\)-\(\Oh(1)\) bằng chia khối. Chia mảng thành các
khối độ dài \(L\) xấp xỉ \(\frac12\log_2 n\). Với mỗi khối, tiền xử lý mọi
truy vấn nội khối; số mẫu hình nội khối của dãy độ sâu Euler bị ràng buộc
bởi bước \(\pm1\), nên số kiểu là \(2^{L-1}\), và tổng bảng kiểu vẫn nhỏ hơn
tuyến tính khi chọn \(L\) như trên.

Giữa các khối, ta lấy giá trị nhỏ nhất của từng khối rồi xây sparse table
trên dãy khối. Một truy vấn bất kỳ tách thành đuôi khối trái, các khối trọn
ở giữa và đầu khối phải. Hai phần biên dùng bảng nội khối, phần giữa dùng
sparse table cấp khối. Với dãy Euler của LCA, ràng buộc \(\pm1\) giúp tiền xử
lý toàn bộ đạt \(\Oh(n)\) và truy vấn vẫn \(\Oh(1)\).

\section{Ghi chú cài đặt}

Cần thống nhất RMQ trả về vị trí hay giá trị; với LCA phải trả về vị trí trong
Euler rồi ánh xạ ngược sang đỉnh. Khi có nhiều vị trí cùng độ sâu nhỏ nhất,
chọn vị trí nào cũng cho cùng LCA nếu đoạn Euler được lấy giữa hai lần xuất
hiện đầu tiên. Với sparse table, nên chuẩn bị bảng log để tránh tính
\(\log\) nhiều lần.

Khi dùng Tarjan offline, mọi truy vấn phải được gắn vào danh sách của cả hai
đầu mút. Chỉ trả lời khi đầu kia đã được tô đen, tức cây con của nó đã xử lý
xong; nếu không, truy vấn sẽ được trả lời ở thời điểm đầu kia được duyệt.

\section{Kết luận}

RMQ và LCA là hai bài toán nền có thể chuyển hóa qua lại. Sparse table là lựa
chọn ngắn, dễ cài, rất đáng dùng khi \(\Oh(n\log n)\) bộ nhớ và thời gian chấp
nhận được. Tarjan phù hợp với truy vấn offline. Block RMQ cho lời giải tuyến
tính hơn nhưng cài đặt phức tạp hơn. Trong các bài đồ thị như Kruskal theo
ngưỡng, việc nhận ra cấu trúc cây ẩn thường là bước quyết định.

\end{document}
